% 1.1 引言
\section{引言}

人类对对称性的兴趣可以追溯到许多个世纪之前（Belov 1956\textsuperscript{b}，Coxeter and Moser 1965，Steno 1669），但对它的现代科学研究，通常认为肇始于阿贝 Haüy（Abbé Haüy）。Haüy 研究了方解石标本解理时的行为：通过把它越敲越碎并观察碎片各晶面之间的夹角，他确信晶体是由大量相同单元重复堆砌而成的。Haüy（1815\textsuperscript{a}–\textsuperscript{d}）还研究了其他多种晶体，并把结论总结为其所谓的\textit{对称性定律}（\textit{Loi de symétrie}）。随着点群、布拉维格子（Bravais lattice）与空间群（space group）等观念在十九世纪逐步形成，对称性的研究也随之发展起来。

\textit{点群}（point group）是作用在一个点上、并满足数学意义上成群要求的一组对称操作：晶体学点群还须满足一个附加条件，即必须与某个空间格子相容。在真实晶体中实际出现的对称操作组合只有有限多种。这 32 个点群的推导由 Hessel（1830）发表，但他的工作被忽视了三十多年，直到 Gadolin（1869）才重新导出它们。此后点群得到了广泛研究，既在其原有的晶体学背景下，也在较晚近的、针对分子与固体物理和化学的群论研究背景下。有用的晶体学著作例如 Buerger（1956）与 Phillips（1963\textsuperscript{a}）。与点群群论应用相关的理论综述可参见许多作者（例如 Bhagavantam and Venkatarayudu 1962，Cracknell 1968\textsuperscript{b}，Hamermesh 1962，Heine 1960，Tinkham 1964）。

我们还可以考虑另一类群，即通过平移对称操作得到的群。如果观察晶体的内部结构，会发现它由大量按规则排列的原子或分子组成；可以在晶体内部找到一组彼此"相似"的点——也就是说，从其中任意一点看晶体，与从另一点看完全相同。这样一组相同的点在数学上称为\textit{格子}（lattice）。可以证明：本质不同的、使每个点的环境都相同的排列方式只有少数几种。这项工作由 Bravais（1850）完成，他证明在三维空间中只可能有 14 种不同的格子；因此它们如今被称为\textit{布拉维格子}（Bravais lattices）。尽管 Frankenheim 更早些曾（不正确地）导出过 15 种格子。

点群关心的是有限物体的对称性，天然晶体只有 32 种点群；布拉维格子关心的则是数学点在空间中的排列。要完整研究晶体的内部结构，也就是原子在晶胞中的确切排布，还需要对称性理论的进一步发展，即\textit{空间群}（space group）。空间群考虑的是由若干相同对象构成的排布的对称性，其中每个对象不再是点，而是具有自身对称性的有限物体或原子团。空间群中出现的操作可以是点群中也有的那种操作，即纯转动、反射、反演以及旋转反演（或旋转反射）；但还可能有其他操作：即带平移的螺旋旋转（screw rotation）与滑移反射（glide reflection）——这两类操作中，转轴或普通反射面在转动或反射之外还叠加了晶体整体的平移。关于 230 个空间群的推导，通常认为归功于 Fedorov 与 Schönflies，有时还加上 Barlow。Burckhardt（1967）的文章回顾了空间群推导的历史，并列出了 Barlow、Fedorov 与 Schönflies 的出版物清单。空间群推导的源头是 Jordan（1868, 1869）与 Sohncke（1879）的工作。Sohncke 导出了只含纯转动的 65 个空间群，并指出 Jordan 已在数学上先行导出，只是没有把结果翻译成几何的语言。Schönflies 重新导出了这 65 个空间群，并把理论推广到含反射面的空间群（Schönflies, 1887\textsuperscript{a}, \textsuperscript{b}, 1889, 1891）。Fedorov（1885, 1891\textsuperscript{a}）得到了类似结果，但其著作以俄文写成，在西欧传播不广；关于 Fedorov 生平与工作以及他的出版物清单，可参见 Shafranovskiĭ（1963）的专著（俄文）。可以肯定这两位学者是独立开展工作的——Fedorov 当时是乌拉尔某矿山的负责人，而 Schönflies 是在 Göttingen 的 F.~Klein 建议下开展研究的；不过随着时间推移他们了解了彼此的工作并互相比对过结果。Barlow（1883）最初研究的是球的密堆积，随后同样从 Sohncke 的 65 个群出发，通过加入反射操作得到了其余的空间群（Barlow 1894）。Burckhardt（1967）的结论是：尽管 Schönflies 并非确立 230 个空间群存在性的第一人，他的著述却使全世界科学界普遍了解并确认了这套 enumeration 与 identification。他的工作虽然比 Fedorov 稍晚且完全独立，最终汇聚成书\textit{Krystallsysteme und Krystallstructur}（Schönflies 1891）。Schönflies 写给 Fedorov 的一封信（转引自 Burckhardt 1967）这样说道："我对与你观点一致感到由衷的高兴；我尤其高兴的是，在这一点上我不再孤独；要赢得晶体学家们的认可还需要付出巨大的努力。\textit{我很乐意把优先权让给你}，这对我并非头等大事。"现代记号下方便而详细的空间群清单见《X 射线晶体学国际表》第一卷（\textit{International tables for X-ray crystallography}，Henry and Lonsdale 1965）。目前已鉴定出约 9000 种化合物的空间群（近期清单见 Donnay, Donnay, Cox, Kennard, and King (1963)，Nowacki, Edenharter, and Matsumoto (1967)，以及 Wyckoff (1963, 1964, 1965, 1966, 1968)）。二维空间群的发现已湮没在久远的古代——它们出现在许多文明的墙纸图案与地砖铺装中（例如见 Coxeter and Moser (1965) 第 33 页）。

尽管 X 射线方法已研究过数量庞大的晶体结构，并把晶体归入相应的空间群，但大约在 1890 年完成 230 个空间群的推导之后，对称性理论的研究进展甚微；直到 1951 年 Shubnikov 出版了《有限图形的对称性与反对称性》（俄文版；该书现已译成英文，附有关于 Shubnikov 其他工作的大量文献清单，见 Shubnikov and Belov 1964）。Koptsik（1967\textsuperscript{a}）回顾了对称性理论近 50 年的发展，\textit{Kristallografiya}（《苏联物理学：晶体学》，\textit{Soviet Phys. Crystallogr.} 英译本）第 2 卷卷首刊有 A.~V.~Shubnikov 的传略（1957）。新的进展源于引入\textit{反对称}（anti-symmetry）操作。经典理论——点群、布拉维格子与空间群——本质上都是三维的：点 $P$ 由向量 $\mathbf{r}\ \{=(x, y, z)\}$ 给定，我们考察对称操作对该点的作用。Shubnikov 的基本思想是：除了点的通常坐标 $x$、$y$、$z$ 之外，再给每点附加第四个坐标 $s$，它只能取两个可能值。$s$ 可以是粒子的自旋，两个允许值对应自旋向上与向下；或者纯抽象地，用黑白两种颜色表示。如果引入坐标 $s$ 而各原子的 $s$ 值又是随机指定的，格子的对称性就被完全破坏；但如果自旋都平行于某一方向，或按某种规则排列，对称性就可能保留一部分。引入一个新操作——称之为\textit{反对称操作}（operation of anti-symmetry）$\mathcal{R}$——并与所有普通的点群、空间群操作结合，就能得到一整批新的点群和空间群，称为\textit{黑白群}（black and white groups，或 \textit{Shubnikov 群}）。黑白群的思想其实早于 Shubnikov：Heesch（1929\textsuperscript{a}, \textsuperscript{b}, 1930\textsuperscript{a}, \textsuperscript{b}）已提出，Woods（1935\textsuperscript{a}–\textsuperscript{c}）也讨论过；只是当时看来它们对描述物理系统用处不大。直到中子衍射技术出现，人们才明白这些群可用于描述磁性有序结构。若把 $s$ 取为磁矩相对某方向的平行与反平行两个取值，则 $\mathcal{R}$ 就是使磁矩反转的操作；此时可以把 $\mathcal{R}$ 看作\textit{时间反演}（time-inversion）操作。

有限群理论始于 Cauchy\footnote{1789--1857。}，他注意到许多看似无关的事实可以通过引入群的概念得到统一解释。Galois\footnote{1811--1832。} 为该理论增添了许多新概念，包括不变子群的概念；他方程理论中的部分工作也是群论在应用中威力的第一个惊人例证。不过，第一部系统的群论著作应归功于 Serret（1866）。此后文献数量持续增长，今天抽象群论仍是重要的研究方向。此外，有限群的应用遍及排列理论、对称性研究、代数方程与微分方程等众多数学领域，可以说凡是需要数学方法的学科都离不开群。自然科学中大量问题需要群论知识；可以断言，生物科学、甚至日益数学化的社会科学，也将产生更多有趣的应用。

在数学理论的发展史上，往往能指出若干位推动了主要进展的学者。群论也不例外。唯一让我们犹豫的是，提到某些名字就可能遗漏了其他同样重要的人。但列举 Sylow、Frobenius、Burnside、Schur、Miller 与 Mackey 应无太大争议（同时向 Noether、Brauer、Ito 以及许多其他对抽象群理论有重大贡献、但其工作与本书应用不直接相关的学者致歉）。

Sylow（1872）在描述有限群结构方面取得重大进展，特别是群元素个数可分解为素数乘积时的结构。Frobenius（1896\textsuperscript{a}, \textsuperscript{b}, 1898）开创并主要负责了群表示与群特征标的理论；Burnside（1903, 1911）作了重要的简化并得到许多原创结果，同样必须被视为极有影响的群论学者。Miller 是一位多产的学者（1894 至 1946 年间约 800 篇论文；Miller 1894, 1946），他对各种有限阶群的结构与性质的研究作出了重要贡献：确定了各种特定阶数的有限群的数目，并通过生成关系研究了这些群结构之间的相互关系。

群与其不变子群的表示之间的关系研究不可避免地导向射影表示。Schur 最先注意到这一点，并在一系列出色的决定性论文（1904, 1907, 1911）中不仅奠定了一般射影表示理论的基础，而且建立了大部分至今仍被视为特别重要的结果。又是 Frobenius（1898）首先构造了如今称为诱导表示（induced representation）的对象。不过，这个对物理应用极其重要的概念，直到 1950 年之后才由 Mackey 在一系列论文（1951, 1952, 1953\textsuperscript{a}, \textsuperscript{b}, 1958）中取得极重要的发展，如今不仅用于空间群理论，而且遍及理论物理各领域（另见 Mackey (1968)）。

前面我们描述了点群、布拉维格子与空间群在刻画晶体宏观对称性（测角术）以及内部结构对称性（X 射线或中子衍射实验）方面的重要性。经典物理中也有群论的应用，例如分子或固体简正振动模式的研究（Wigner 1930），或者对属于某点群的晶体，确定其宏观性质张量各分量之间可能存在的关系（例如见 Nye (1957)）。然而，只有随着量子力学的出现，群论与表示理论的全部数学威力才真正成为理解物理体系的最有用工具。将群论应用于量子力学的开创性工作主要出自 Weyl、Wigner 与 von Neumann（见 Weyl (1931)，Wigner 经典著作的英译本 (1959)，以及 von Neumann 文集 (1961, 1963)）。在微观层面研究晶体时必须记住，晶体中的每个粒子都服从量子力学而非经典力学，因此必须用适当的波函数 $\psi$ 来描述。群论用于量子力学的关键结果是 Wigner 经典著作第 11 章阐述的定理（英译本 Wigner (1959)）。若量子体系由适当的 Schrödinger 波动方程描述，则 Wigner 定理可概括为："\textit{属于某一特定本征值的 Schrödinger 方程群的表示，在相似变换的意义下唯一确定}。"除偶然简并外，该表示是不可约的。不可约表示之所以重要，是因为它们可以用来无歧义地标记量子体系的能级。晶体学点群与双值点群的不可约表示早在 Bethe（1929）就已确定，并广泛用于分子能级的标记（综述与专著见 Eyring, Walter, and Kimball (1944)，Nussbaum (1968)，Rosenthal and Murphy (1936)，Slater (1963)，Wilson, Decius, and Cross (1955)），晶体中离子或分子能级的标记（见 Herzfeld and Meijer (1961)，Hutchings (1964)，Judd (1963)，McClure (1959\textsuperscript{a}, \textsuperscript{b})），以及晶体激子的标记（Overhauser 1956）。Koster, Dimmock, Wheeler, and Statz（1963）对晶体学点群及其表示的重要性质作了特别有用的总结。

空间群不可约表示的理论由 Seitz（1936\textsuperscript{b}）建立，最初由 Bouckaert, Smoluchowski, and Wigner（1936）应用于点式空间群，Herring（1942）应用于非点式空间群，Elliott（1954）应用于双值空间群。此后许多作者确定了各个空间群的不可约表示，并采用了多种不同的记号。Koster（1957）写了重要的综述；近来还出版了全部 230 个空间群不可约表示的完整表格（Faddeev 1964，Kovalev 1965，Miller and Love 1967，Zak, Casher, Glück and Gur 1969）。空间群不可约表示的重要性在于：根据 Wigner 定理，它们可用于标记晶体中粒子或准粒子的能级，从而标记晶体的电子能带结构与声子色散曲线（电子能带综述见 Blount (1962)，Jones (1960)，Nussbaum (1966)，Slater (1965\textsuperscript{b}, 1967)；声子色散曲线见 Maradudin and Vosko (1968)，Warren (1968)）。类似地，空间群的不可约表示也可用于标记磁性晶体中的磁振子色散曲线。不过这里有一个额外的复杂之处：黑白 Shubnikov 空间群拥有的是\textit{共表示}（corepresentations）而非普通表示（Dimmock and Wheeler 1962\textsuperscript{b}，Karavaev, Kudryavtseva, and Chaldyshev 1962，Loudon 1968，Wigner 1959, 1960\textsuperscript{a}, \textsuperscript{b}）。

如果确定点群与空间群不可约表示的全部心血仅仅换来一套能级标记方案，那么这些努力是否值得就值得怀疑了。然而，不可约表示还能用来确切判定波函数 $\psi_{i}$ 在分子或晶体的 Schrödinger 群的各个操作下如何变换。这使得在把未知波函数按已知函数（如球谐函数）展开时可以进行重要的简化（Altmann 1957，Altmann and Bradley 1963\textsuperscript{b}，Bell 1954，Betts 1959，McIntosh 1963，von der Lage and Bethe 1947），或按平面波展开（Cornwell 1969，Luehrmann 1968，Schlosser 1962，Slater 1965\textsuperscript{b}, 1967）。把未知函数 $\psi_{i}$ 在能级 $E_{i}$ 的展开限制在已知属于 $E_{i}$ 的表示的那些函数上，往往能大大简化求解 Schrödinger 方程以确定 $\psi_{i}$ 的实际计算。当体系在两个能级 $E_{i}$ 与 $E_{j}$ 之间跃迁时，波函数 $\psi_{i}$ 的变换性质同样重要：利用跃迁的量子力学矩阵元必须是纯数这一条件，可以对任意给定的微扰势判断跃迁是允许还是禁戒，即确定\textit{选择定则}（selection rules）。孤立分子以及晶体中的离子或分子的跃迁选择定则的群论判定，涉及各种点群表示乘积的研究，相关讨论见前文已引的文献。利用 $\psi_{i}$ 的变换性质研究晶体中非局域态的跃迁选择定则则更为复杂，初步工作已由若干作者完成（Birman 1962\textsuperscript{b}, 1963，Elliott and Loudon 1960，Lax and Hopfield 1961，Zak 1962）。
